
\section*{2. Premières réductions}

\textbf{2.1.} Appliquons 1.3 à \(X/S\) et \(Y/S\)~: on trouve un ouvert non vide \(V\) de \(S\), un schéma intègre \(S'\) et un morphisme fini plat \(S'\to V\) tels que toute composante irréductible \(X'_i\) de \(X_{S'}\) et \(Y'_j\) de \(Y_{S'}\), munie de la structure réduite induite, soit intègre sur \(S'\). En vertu de la descente f.p.q.c. (EGA IV 2.7.1 (xiv) et SGA 1 VIII 7.9), on peut remplacer \(S\) par \(S'\). Il résulte alors du diagramme commutatif suivant~:
\[
\begin{tikzcd}
X_{\red} \arrow[r,"f_{\red}"] \arrow[d,"i"] & Y_{\red} \arrow[d,"j"]\\
X \arrow[r,"f"] & Y & ,
\end{tikzcd}
\]
impliquant le diagramme commutatif
\[
\begin{tikzcd}
\underline{\Pic}_{X_{\red}/S} & \underline{\Pic}_{Y_{\red}/S} \arrow[l]\\
\underline{\Pic}_{X/S} \arrow[u] & \underline{\Pic}_{Y/S} \arrow[l] \arrow[u] & ,
\end{tikzcd}
\]
et de sorites triviaux sur les morphismes de foncteurs représentables par morphismes quasi-affines (résultant de EGA II 5.1.10 (ii) (v)) qu'il suffit de démontrer deux assertions suivantes :
\begin{enumerate}[label=\Roman*.]
\item \uline{Le théorème} 1.1 \uline{lorsque} \(f\) \uline{est une immersion (nécessairement nilpotente).}
\item \uline{Le théorème} 1.1 \uline{lorsque} \(X\) \uline{et} \(Y\) \uline{sont réduits et que toute composante irréductible} \(X_i\) \uline{de} \(X\) \uline{et} \(Y_i\) \uline{de} \(Y\)\uline{, munie de la structure réduite induite, est intègre sur} \(S\).
\end{enumerate}

\textbf{2.2.} Étudions l'assertion II. Pour tout \(i\), soit \(j=\varphi(i)\) un indice tel que \(f(X_i)\subset Y_j\), soit \(\Sigma_j=\coprod_{\varphi(i)=j} X_i\). Alors on a les deux diagrammes
